
\documentclass{article}

%--- Document Formatting ---%
\usepackage{lmodern} % gives you the more modern Latin Modern fonts
\usepackage[T1]{fontenc} % fontenc provides better and more character encoding
\usepackage[margin=3cm]{geometry} % geometry allows changing page sizes, orientation and margin
\usepackage{microtype}

\usepackage{titlesec} % titlesec allows changing the style, spacing and font of chapter, sections and subsections
% {left-indent}{space before}{space after}
\titlespacing*{\section}{0pt}{6ex plus 1ex minus .2ex}{2.3ex plus .2ex}
\titlespacing*{\subsection}{0pt}{2.5em plus 1ex minus .2ex}{2ex plus .2ex}
\linespread{1.2}


%--- Language ---%
\usepackage[australian]{babel}
\usepackage[useregional]{datetime2}


%--- Symbols ---%
\usepackage{amssymb} % amssymb provides symbols, mainly number systems
\usepackage{cancel} % adds \cancel{x}, allowing striking out derivations
%\usepackage{siunitx} % used for physics based units


%--- Tools and General Formatting ---%
\usepackage{mathtools} % an extension for amsmath, loads amsmath itself
\usepackage{amsthm} % improves proof based math
\usepackage{enumitem} % enumitem gives more control over lists

%\usepackage{graphicx} % use to include images/figures
%\usepackage{xcolor} % for colour
%\usepackage[hidelinks]{hyperref} % adds hyperlinks, can sometimes cause issues. Keep last
%\usepackage{cleveref} % extension for hyperref, auto formatted cross-reference


\addtolength{\jot}{1em} % Increases row separation locally

\begin{document}

\title{Algebra Revision}
\author{Montague Hamilton}
\date{\today}
\maketitle

\section{Audit}
\subsection{Indices}
\begin{enumerate}
  \item Simplify:
    \begin{enumerate}[label=\alph*)]
    \item \begin{align*}
        2^{5} \cdot 2^{-3} \div 2^{4} &= 2^{5-3-4} \\
                      &= 2^{-2}
      \end{align*}
    \item \begin{align*}
        (x^{3})^{2} \cdot x^{-4} &= x^{3 \cdot 2}\cdot x^{-4} \\
                    &= x^{6 - 4} \\
                    &= x^{2}
      \end{align*}
  \end{enumerate}
  \item Evaluate, no calculator:
    \begin{enumerate}[label=\alph*)]
      \item $27^{\frac{2}{3}} = \sqrt[3]{27^{2}} = \sqrt[3]{729} = 9$
      
      \item \begin{align*}
          \left(\frac{1}{16} \right)^{\frac{-3}{4}} &= \left(\frac{16}{1} \right)^{\frac{3}{4}} \\
          &= \left(2^{4} \right)^{\frac{3}{4}} = 2^{3} \\
          &= 8 % Took quite a while with this one, forgot the underlying rule, not tied to understanding
        \end{align*}

    \item $4^{\frac{-1}{2}} = \sqrt[2]{4^{-1}} = \sqrt[2]{\frac{1}{4}} = \frac{1}{2}$
    \end{enumerate}

  \item Expand and simplify:
  \begin{enumerate}[label=\alph*)]
    \item \begin{align*}
        (x + 3)(x - 5) &= x^{2} - 5x + 3x - 15 \\
                        &= x^{2} -2x -15
      \end{align*}

    \item \begin{align*}
        (2x - 1)^{2} &= (2x - 1)(2x - 1) \\
                  &= 4x^{2} - 2x - 2x + 1 \\
                  &= 4x^{2} - 4x + 1
      \end{align*}
  \end{enumerate}

  \item Factorize:
    \begin{enumerate}[label=\alph*)]
      \item $x^{2} - 9 = (x + 3)(x - 3)$
      \item $x^{2} + 5x + 6 = (x + 2)(x + 3)$
      \item $6x^{2} + x - 2 = (3x + 2)(2x - 1)$
    \end{enumerate}

  \item Factorize Completely: 
    \begin{align*}
      x^{3} - x &= x(x^{2} - 1) \\
      &= x((x + 1)(x - 1))
    \end{align*}
    


  \item Solve, giving all solutions:
    \begin{align*}
      x^{2} &= 5x \\
      x  &= \pm5
    \end{align*}
    Attempt 2:
    \begin{align*}
      2x^{2} &= 5x \\
      2x^{2} - 5x &= 0 \\
      x(2x - 5) &= 0 
    \end{align*}
    

    \item Solve:
    \begin{align*}
      x^{2} + 2x - 8 &= 0 \\
      (x + 4)(x - 2) &= 0
    \end{align*}
    The quadratic has x intercepts at $(0, -2)$ and $(0, 4)$
    Attempt 2:
    Assuming the intercepts are at $(-2, 0)$ and $(4, 0)$
    \begin{align*}
      (-2)^{2} + 2(-2) -8 &= 0 \\
      4 - 4 - 8 &\neq 0
    \end{align*}
    My factorize term under expansion:
    \begin{align*}
      (x + 4)(x - 2) &= 0 \\
      x^{2} - 2x + 4x - 8 &= 0 \\
      x^{2} + 2x - 8 &= 0 
    \end{align*}
    So the factorized term is correct, at second glance, coordinates should be $(-4, 0)$ and $(2, 0)$, testing:
    \begin{align*}
      (-4)^{2} + 2(-4) - 8 &= 0 \\
      16 - 8 - 8 &= 0 \\
      0 &= 0
    \end{align*}
    So that one works, testing the next one:
    \begin{align*}
      (2)^{2} + 2(2) - 8 &= 0
      4 + 4 - 8 &= 0 \\
      0 &= 0
    \end{align*}
    Finding the y intercept:
    \begin{align*}
      y &= (0)^{2} + 2(0) - 8 \\
      y &= 0 + 0 - 8 \\
      y &= -8
    \end{align*}
    So the y intercept is at $(0, -8)$
    
    
  \item Make $h$ the subject:
    \begin{align*}
       V &= \frac{1}{3} \pi \cdot r^{2} \cdot h \\
      h &= \frac{3 V}{\pi \cdot r^{2}}
    \end{align*}

  \item Solve simultaneously:
    \begin{equation}
      x + y = 7
    \end{equation}
    \begin{equation}
      2x - 3y = 4
    \end{equation}
    \begin{align*}
      (1) x + y &= 7 \\
      x &= -y + 7 \\
      y &= -x + 7
    \end{align*}
    Sub (1) into (2), first solving $y$, then $x$: 
    \begin{align*}
      2(-y + 7) - 3y &= 4 \\
      -2y + 14 -3y &= 4 \\
      -5y &= -10 \\
      y &= 2
    \end{align*}
    \begin{align*}
      2x - 3(-x + 7) &= 4 \\
      2x + 3x - 21 &= 4 \\
      5x &= 25 \\
      x &= 5
    \end{align*}
    We can then confirm, first (1) then (2):
    \begin{align*}
      (5) + (2) &= 7 \\
       7 &= 7 \\
      \text{and} \\
      2(5) - 3(2) &= 4 \\
      10 - 6 &= 4 \\
      4 &= 4
    \end{align*}

  \item Simplify: $\frac{x^{2} - 4}{x -2} = x - 2$ where $x \neq 2$
    Attempt 2: First factor $x^{2} - 4$:
    \begin{equation}
      x^{2} - 4 = (x + 2)(x - 2)
    \end{equation}
    Then sub into original problem:
    \begin{align*}
      && \frac{x^{2} - 4}{x - 2} &= \frac{(x + 2)(x - 2)}{x - 2} \\
      && \frac{(x + 2)\cancel{(x - 2)}}{\cancel{x - 2}} &= x + 2 && \text{where } x \neq 2 
    \end{align*}
    
    

  \item Express as a single fraction: $\frac{1}{x} + \frac{1}{x + 1} = \frac{2}{2x + 1}$ \\
    Attempt 2: $x \cdot (x + 1) = x^{2} + x$
    \begin{align*}
      \frac{1}{x} + \frac{1}{x + 1} &= \frac{x + 1}{x^{2} + x} + \frac{x}{x^{2} + x} \\
      &= \frac{2x + 1}{x^{2} + x}
    \end{align*}
    

  \item Rationalize the denominator: $\frac{2}{\sqrt[2]{3} - 1} = \frac{4}{3 + 1} = 1$ \\
    Attempt 2, focusing on the denominator first: 
    \begin{align*}
      && (\sqrt[2]{3} - 1)^{2} &= 3 - \sqrt[2]{3} - \sqrt[2]{3} + 1 && \text{not the one} \\
      && (\sqrt[2]{3} - 1)(\sqrt[2]{3} + 1) &= 3 + \sqrt[2]{3} - \sqrt[2]{3} - 1 && \text{this is the one}
    \end{align*}
    Now for the actual equation:
    \begin{align*}
      \frac{2}{\sqrt[2]{3} - 1} &= \frac{2}{\sqrt[2]{3} - 1} \cdot \frac{\sqrt[2]{3} + 1}{\sqrt[2]{3} + 1} \\
      \text{sub from above} &= \frac{2(\sqrt[2]{3} + 1)}{2} \\
      \frac{\cancel{2}(\sqrt[2]{3} + 1)}{\cancel{2}} &= \sqrt[2]{3} + 1
    \end{align*}
    checking, subbing $\sqrt[2]{3}$ for 1.73:
    \begin{align*}
      \frac{2}{1.73 - 1} &\approx 2.74 \\
      1.73 + 1 &\approx 2.73 \approx 2.74
    \end{align*}
    
    

\end{enumerate}
\end{document}
